\section{Calibration of the finite certificates}\label{app:certificates}

We verify the central-interval routine and choose the finite calculation's
parameters. The purpose of the calibration is to ensure that a failed
certificate identifies a coordinate whose short interval makes a
middle-half refinement profitable.

\subsection{The central-interval test}

For the one-coordinate list $[A_i,B_i]$, $1\le i\le N$, take $u$ to be
the $N/4$-th upper endpoint in sorted order and $v$ to be the $3N/4$-th
lower endpoint. The alphabet size $S$ is divisible by four, and the queried
coordinate is interior, so $N/4$ is an integer. Return \texttt{FAIL} if
$u\ge v$ or fewer than $N/4$ list intervals lie wholly in $[u,v]$.
Otherwise return $[u,v]$.

\begin{proof}[Proof of Lemma~\ref{lem:central-interval}]
At least $N/4$ intervals end at or below $u$, and at least $N/4$ begin
at or above $v$. On success, at least $N/4$ intervals also lie wholly
inside $[u,v]$. Enclosure certifies mass at least $1/4$ in each of the
three regions. The true distribution is continuous and supported on
$(\ell_x,r_x)$, so these masses imply $\ell_x<u<v<r_x$.
The endpoints are dyadic, and sorting and counting finite lists are
deterministic.
\end{proof}

\subsection{The width at which certificates succeed}

The graph bounds limit how many coordinates can be comparable with a
queried coordinate. Define
\begin{equation}\label{eq:comparable-bound}
B=(k+1)^{2H+1}.
\end{equation}
In an in-node DAG of degree at most $k$ and directed paths of at most
$2H$ vertices, every comparable coordinate lies within undirected distance
$2H$. Thus $B$ bounds the number of other comparable coordinates.
Set
\begin{equation}\label{eq:width-threshold}
w_*=2^{-2kB}.
\end{equation}
This width calibrates the finite calculations. The algorithms test the
computed certificates directly.

\begin{lemma}[When certification succeeds]\label{lem:certificate-success}
Suppose a node has positive feasible volume, its in-node DAG has degree
at most $k$, and its directed paths have at most $2H$ vertices.
If any free coordinate involved in a comparison has interval width at
least $w_*$, comparison certification succeeds. If $W_x\ge w_*$, then
\textnormal{\texttt{CentralInterval}}$(v,x)$ succeeds.
\end{lemma}

\begin{proof}[Proof of Lemma~\ref{lem:certificate-success}]
Calibrate both routines using
Theorem~\ref{thm:local} with
\begin{equation}\label{eq:estimator-coupling-parameters}
d=k,\qquad h=2H,\qquad
\alpha=\frac\eta2,\qquad
\delta=\frac{\eta w_*}{8(B+1)}.
\end{equation}
These parameters are fixed for the interaction. They do not depend on
the queried widths.

Suppose one queried free coordinate $x$ has $W_x\ge w_*$.
If the computed intervals overlap and both have width at most $\delta$,
the enclosed positions differ by at most $2\delta$.
Lemma~\ref{lem:gap} therefore bounds the unresolved fraction by
\[
1-p-q\le\alpha+\Prb(|X_x-X_y|\le2\delta)
\le\alpha+\frac{4(B+1)\delta}{w_*}=\eta.
\]
Thus $\max(p,q)\ge(1-\eta)/2>1/2-\eta$, so comparison certification
succeeds. For a threshold $c$, an unresolved interval contains $c$.
If its width is at most $\delta$, its enclosed position is within
$\delta$ of $c$, giving the stronger bound
\begin{equation}\label{eq:threshold-unresolved}
\alpha+\frac{2(B+1)\delta}{w_*}=\frac{3\eta}{4}.
\end{equation}

It remains to show that \texttt{CentralInterval} succeeds when
$W_x\ge w_*$. By \eqref{eq:threshold-unresolved}, at most a fraction
$\tau=3\eta/4$ of its intervals contain any fixed real number.
This bound applies in particular to the selected endpoints $u$ and $v$.
An interval extending below $u$ either ends strictly below $u$ or
contains $u$. By the rank defining $u$, fewer than $N/4$ intervals end
strictly below it. Thus at most $(1/4+\tau)N$ intervals extend below $u$.
Similarly, at most $N/4$ intervals begin strictly above $v$, so at most
$(1/4+\tau)N$ extend above $v$.

At least $(1/2-2\tau)N\ge N/4$ intervals therefore lie wholly inside
$[u,v]$, since $\eta<1/32$. Enclosure gives this interval mass at least
$1/4$. Continuity forces $u<v$, so both success tests pass.
\end{proof}
